\documentclass[aspectratio=169,11pt]{beamer}
\usepackage[T1]{fontenc}
\usepackage{lmodern,booktabs,tabularx,array}
\definecolor{navy}{HTML}{16324F}
\definecolor{muted}{HTML}{536272}
\setbeamercolor{structure}{fg=navy}
\setbeamercolor{frametitle}{fg=navy}
\setbeamerfont{frametitle}{size=\Large,series=\bfseries}
\setbeamertemplate{navigation symbols}{}
\setbeamersize{text margin left=7mm,text margin right=7mm}
\setbeamertemplate{footline}{\leavevmode\hbox{\begin{beamercolorbox}[wd=.85\paperwidth,ht=2.6ex,dp=1.1ex,leftskip=5mm]{author in head/foot}\color{muted}\scriptsize OPALX \enspace|\enspace KV tune-shift benchmark\end{beamercolorbox}\begin{beamercolorbox}[wd=.15\paperwidth,ht=2.6ex,dp=1.1ex,rightskip=5mm]{date in head/foot}\hfill\color{muted}\scriptsize\insertframenumber/\inserttotalframenumber\end{beamercolorbox}}}
\newcolumntype{Y}{>{\raggedright\arraybackslash}X}
\newcommand{\src}[1]{\par\vfill{\color{muted}\scriptsize #1\par}}
\newcommand{\epsunit}{\ensuremath{\pi\,\mathrm{mm\,mrad}}}
\renewcommand{\arraystretch}{1.15}
\title{OPALX KV tune-shift benchmark}
\author{OPALX}
\date{25 November 2026}
\begin{document}
\begin{frame}{OPALX: benchmark parameters}
\small
\begin{tabularx}{\linewidth}{lY}
\toprule Parameter & Value used \\
\midrule
Lattice and reference radius & 16 FD cells; $R_0=12.0$ m \\
Field index and spiral angle & $k=7.630003$; $49.973388^\circ$ \\
Nominal F / D fields & $+1.650$ T / $-0.7673721$ T \\
F / D widths; F--D gap & $5.9083^\circ$ / $2.95415^\circ$; $2.38755^\circ$ \\
Fringe & Product of sigmoids; $W_0=1.17W$, steepness $16/W_0$ \\
Beam and population & Protons, $300$ MeV; $2.08\times10^{13}$ protons/bunch \\
Transverse beam & KV; equal-plane emittance cases $50,70,90$ \epsunit \\
Timestep & $0.02$ ns \\
Mesh & $128\times128$ transverse; 763 full-ring longitudinal cells \\
Longitudinal resolution & $127.9945$ cells across the initial 12 m bunch \\
\bottomrule
\end{tabularx}\par
\vspace{3mm}
\footnotesize Table I supplies the machine values, timestep and target mesh counts. Exact lattice and fringe values come from the public PyFFA configuration.
\src{Zhou et al., Phys. Rev. Accel. Beams \textbf{29}, 113301 (2026), Table I and Fig. 6; \href{https://doi.org/10.1103/q1xb-z29j}{doi:10.1103/q1xb-z29j}. PyFFA revision \texttt{2cca56b5}.}
\end{frame}

\begin{frame}{Input that we needed to infer}
\small
\begin{tabularx}{\linewidth}{lY}
\toprule Quantity & OPALX choice \\
\midrule
Bunch length and cut & Gaussian $\sigma_s=1.5$ m; truncate at $\pm4\sigma_s$ \\
Sampling & $131\,072$ macroparticles \\
Initial energy spread & Zero \\
Transverse mesh extent & ($0.4\times0.4$ m$^2$)  {\color{red} follow beam}\\
Magnetic expansion & (2) {\color{red} 5} order \\
Turns for tune calculation & (10)  {\color{red} 200}  \\ 
Emittance convention & Numeric rms $\epsilon=(50,70,90)\times10^{-6}$ m rad; no extra $\pi$ \\
\bottomrule
\end{tabularx}\par
\vspace{3mm}
\footnotesize These are guessed settings, {\color{red} red} now resolved. 
\end{frame}

%\begin{frame}{OPALX tunes and tune shifts over fit windows}
%The following tables report central-bunch medians. A median is not the maximum-depression endpoint used to describe the theoretical core.
%\small
%\begin{center}\setlength{\tabcolsep}{7pt}\renewcommand{\arraystretch}{1.08}\input{generated/short-windows.tex}\end{center}
%\vspace{2mm}
%\footnotesize $\Delta Q_i=Q_i-Q_{i0}$, with $(Q_{x0},Q_{y0})=(3.2331550664,2.8132163500)$. Vertical $y$ corresponds to the paper's $z$.
%\src{Medians from eight probe samples per particle revolution. Fit windows start two revolutions after the first crossing and retain the fixed per-axis core cohort. All windows reuse the same tracking histories.}
%\end{frame}

\begin{frame}{Comparison with the paper}
\small
\begin{center}\setlength{\tabcolsep}{5pt}\input{generated/short-paper.tex}\end{center}


\begin{itemize}
\item OPALX uses eight-turn medians of accepted central-bunch tunes ($|s_{\rm initial}|<0.15$ m).
\item This is our analysis choice: medians reduce sensitivity to extreme fitted tunes.
\item the paper uses the rounded bare horizontal tune $3.233$.
\item its numerical vertical tune values are not specified in the paper.
\item The paper does not specify how its PIC core tunes were extracted. We cannot establish that it used medians or the same statistic.
\end{itemize}
\end{frame}

\begin{frame}{OPALX KV footprint: $50$ \epsunit}
\centering
\includegraphics[width=\linewidth,height=.79\textheight,keepaspectratio]{figures/opalx-50-footprint-hull.pdf}
\src{Left: lowest accepted tune in each 0.25 m initial slice (maximum depression); sensitive to extreme fitted tunes. Dashed curve: paper's theoretical peak scaled by the adopted Gaussian profile, not digitized paper data. Right: OPALX density with approximate visible-color convex hull from Zhou et al., Fig. 6 (CC BY 4.0); magenta dashed, no fitting or rescaling. Blue circle: bare tune.}
\end{frame}

\begin{frame}{OPALX KV footprint: $70$ \epsunit}
\centering
\includegraphics[width=\linewidth,height=.79\textheight,keepaspectratio]{figures/opalx-70-footprint-hull.pdf}
\src{Left: lowest accepted tune in each 0.25 m initial slice (maximum depression); sensitive to extreme fitted tunes. Dashed curve: paper's theoretical peak scaled by the adopted Gaussian profile, not digitized paper data. Right: OPALX density with approximate visible-color convex hull from Zhou et al., Fig. 6 (CC BY 4.0); magenta dashed, no fitting or rescaling. Blue circle: bare tune.}
\end{frame}

\begin{frame}{OPALX KV footprint: $90$ \epsunit}
\centering
\includegraphics[width=\linewidth,height=.79\textheight,keepaspectratio]{figures/opalx-90-footprint-hull.pdf}
\src{Left: lowest accepted tune in each 0.25 m initial slice (maximum depression); sensitive to extreme fitted tunes. Dashed curve: paper's theoretical peak scaled by the adopted Gaussian profile, not digitized paper data. Right: OPALX density with approximate visible-color convex hull from Zhou et al., Fig. 6 (CC BY 4.0); magenta dashed, no fitting or rescaling. Blue circle: bare tune.}
\end{frame}
\end{document}
